\documentclass[aps,prl,twocolumn,superscriptaddress,showpacs]{revtex4-2}
\usepackage{amsmath,amssymb,amsfonts}
\usepackage{graphicx}
\usepackage{hyperref}

\begin{document}

\title{Letter 73: BCT Signatures in Gravitational Waves\\
\large A Scalar Breathing Mode and Topological Memory\\
Predicted for the Einstein Telescope}

\author{Michel Robert Cabri\'{e}}
\affiliation{Independent Researcher, Victoria, Australia}
\email{ZeroFreeParameters@gmail.com}

\begin{abstract}
The BCT Superfluid Lattice Model makes two falsifiable predictions 
for gravitational wave astronomy beyond general relativity. First, 
a scalar breathing polarisation mode with amplitude 
$h_\mathrm{scalar}/h_\mathrm{tensor} = \alpha_0/v_\mathrm{ph}^2 
= 1.61\times10^{-4}$ arises from isotropic compression of the 
Octet-Hopfion Condensate (OHC). Second, gravitational waves imprint 
a permanent \emph{topological memory} in the OHC lattice — a Hopf 
charge gradient additional to the standard GW memory — with amplitude 
$h_\mathrm{BCT} = \alpha_0 \times h_\mathrm{standard} \approx 
1.30\times10^{-25}$ for GW150914-like events at 410\,Mpc. This places 
the BCT topological memory signal within the design sensitivity of the 
Einstein Telescope ($\sim10^{-25}$). All current gravitational wave 
observations are fully consistent with BCT. The model predicts 
$v_\mathrm{GW} = c$ exactly, no dispersion, and standard GR waveforms, 
in agreement with observations including GW170817.
\end{abstract}

\maketitle

\section{Introduction}

The first direct detection of gravitational waves~\cite{GW150914} 
confirmed one of general relativity's most dramatic predictions. 
Over 300 events have now been catalogued by the LIGO-Virgo-KAGRA 
network~\cite{GWTC3}, and the most recent event GW250114 provides 
the clearest test of GR to date~\cite{GW250114}. All observations 
are consistent with standard GR.

The BCT Superfluid Lattice Model~\cite{BCT_Monograph} derives 
general relativity from the long-wavelength dynamics of the 
Octet-Hopfion Condensate (OHC) vacuum~\cite{BCT_AppJH}. 
Gravitational waves are collective phonon modes of the OHC 
lattice propagating at exactly $c$. In this framework, all GR 
predictions are recovered to leading order, but two sub-leading 
BCT-specific signatures arise that are in principle observable 
with next-generation detectors.

\section{Gravitational Waves in BCT}

In the BCT framework, the spacetime metric is an emergent 
long-wavelength description of OHC condensate dynamics. 
The exterior condensate $\Psi_1$ propagates at speed $c$ 
(the Goldstone mode), recovering special relativity. 
Gravitational waves are transverse oscillations of this 
condensate, identical to GR tensor modes at leading order.

The BCT dispersion relation for gravitational waves is:
\begin{equation}
\omega^2 = c^2 k^2 \left[1 + \alpha_0 (k\,l_P)^2 + 
\mathcal{O}(\alpha_0^2)\right],
\end{equation}
where $\alpha_0 = r_\mathrm{oct}\times r_\mathrm{tet}/\pi = 
0.0074081$ and $l_P = 1.616\times10^{-35}$\,m. The correction 
term $\alpha_0(kl_P)^2 \lesssim 10^{-75}$ at LIGO frequencies 
--- completely negligible. BCT predicts $v_\mathrm{GW} = c$ 
to any measurable precision, consistent with the GW170817 
constraint $|v_\mathrm{GW}-c|/c < 10^{-15}$~\cite{GW170817}.

\section{Prediction 1: BCT Scalar Breathing Mode}

Beyond the two standard GR tensor polarisations ($h_+$ and 
$h_\times$), BCT predicts a third polarisation mode: an 
isotropic scalar ``breathing'' mode arising from compression 
of the OHC condensate along the radial direction.

The physical origin is the interior condensate $\Psi_2$, which 
fills each BCT Planck sphere with Hopf charge $H=1$ and interior 
phase velocity $v_\mathrm{ph} = 6.78\,c$~\cite{BCT_AppJH}. 
A passing gravitational wave deforms the exterior condensate 
$\Psi_1$, which couples through the $\alpha_0$ Josephson junction 
to the interior $\Psi_2$. The mismatch in phase velocities 
($c$ exterior vs $6.78c$ interior) generates a scalar 
breathing response with amplitude:
\begin{equation}
\frac{h_\mathrm{scalar}}{h_\mathrm{tensor}} = 
\frac{\alpha_0}{(v_\mathrm{ph}/c)^2} = 
\frac{0.0074081}{(6.78)^2} = 1.61\times10^{-4}.
\end{equation}

For a GW150914-like event with $h_\mathrm{tensor}\sim10^{-21}$:
\begin{equation}
h_\mathrm{scalar} \approx 1.61\times10^{-25}.
\end{equation}

\begin{table}[h]
\centering
\caption{BCT scalar mode vs detector sensitivities}
\begin{tabular}{ll}
\hline\hline
Detector & Strain sensitivity \\
\hline
Advanced LIGO & $\sim 10^{-23}$ \\
Einstein Telescope & $\sim 10^{-25}$ \\
LISA (massive BH) & $\sim 10^{-22}$ \\
BCT scalar (GW150914) & $1.61\times10^{-25}$ \\
\hline\hline
\end{tabular}
\end{table}

The BCT scalar mode amplitude is at the Einstein Telescope 
design sensitivity for strong nearby events. A non-detection 
at this level would constrain $\alpha_0/(v_\mathrm{ph}/c)^2$ 
and provide a direct test of the BCT interior phase velocity.

\section{Prediction 2: BCT Topological Memory}

The standard gravitational wave memory effect~\cite{GWMemory} 
is a permanent displacement of spacetime following a GW burst. 
BCT predicts an additional \emph{topological memory} component: 
a permanent Hopf charge gradient in the OHC lattice left after 
a gravitational wave passes.

\subsection{Physical Mechanism}

As a gravitational wave deforms the OHC lattice, the Josephson 
coupling $\alpha_0$ drives a small fraction of OHC spheres 
through transient $H=1 \to H=0$ phase slips. After the wave 
passes, the lattice relaxes — but topology enforces that the 
net Hopf charge displaced is conserved ($H\in\mathbb{Z}$), 
leaving a permanent spatial gradient. This is the BCT topological 
scar of the gravitational wave.

\subsection{Calculation}

The standard GW memory strain for a source radiating energy 
$\Delta E_\mathrm{GW}$ at distance $r$ is:
\begin{equation}
h_\mathrm{standard} = \frac{G}{c^4}\frac{\Delta E_\mathrm{GW}}{r}.
\end{equation}

For GW150914 ($\Delta E_\mathrm{GW} \approx 3\,M_\odot c^2$, 
$r = 410$\,Mpc):
\begin{equation}
h_\mathrm{standard}^\mathrm{GW150914} \approx 1.75\times10^{-23}.
\end{equation}

The BCT topological memory is a fixed fraction $\alpha_0$ 
of the standard memory, arising because each OHC Josephson 
junction converts fraction $\alpha_0$ of the passing strain 
energy into topological Hopf charge displacement:
\begin{equation}
\boxed{h_\mathrm{BCT} = \alpha_0 \times h_\mathrm{standard} 
= 1.30\times10^{-25}.}
\end{equation}

\begin{table}[h]
\centering
\caption{BCT topological memory for GW150914}
\begin{tabular}{lll}
\hline\hline
Quantity & Value & Note \\
\hline
$h_\mathrm{standard}$ & $1.75\times10^{-23}$ & Standard GR \\
$h_\mathrm{BCT}$ & $1.30\times10^{-25}$ & BCT addition \\
$h_\mathrm{BCT}/h_\mathrm{standard}$ & $\alpha_0 = 0.74\%$ & Exact \\
ET sensitivity & $\sim10^{-25}$ & Design goal \\
\hline\hline
\end{tabular}
\end{table}

The ratio $h_\mathrm{BCT}/h_\mathrm{standard} = \alpha_0$ is 
fixed by BCT void geometry with zero free parameters:
\begin{equation}
\alpha_0 = \frac{r_\mathrm{oct}\times r_\mathrm{tet}}{\pi} 
= \frac{(\sqrt{2}-1)/2 \times (\sqrt{6}-2)/4}{\pi} = 0.0074081.
\end{equation}

The total GW memory observable is:
\begin{equation}
h_\mathrm{total} = h_\mathrm{standard}(1 + \alpha_0) 
= h_\mathrm{standard}\times 1.007408.
\end{equation}

\section{Consistency with All Current Observations}

BCT is fully consistent with all gravitational wave observations:

\begin{itemize}
\item $v_\mathrm{GW} = c$ exactly (dispersion correction 
$\sim10^{-75}$) ✓
\item Standard tensor polarisations $h_+$, $h_\times$ ✓  
\item GR waveform shape at leading order ✓
\item GW170817 speed constraint $|v-c|/c < 10^{-15}$ ✓
\item GW250114 black hole spectroscopy consistent with GR ✓
\end{itemize}

The two BCT predictions — scalar mode and topological memory 
— are sub-leading effects invisible to current detectors but 
within reach of the Einstein Telescope.

\section{Discussion}

The BCT topological memory has a qualitatively different 
character from the standard GW memory. Standard memory is 
a permanent \emph{geometric} displacement of spacetime. 
BCT topological memory is a permanent \emph{topological} 
displacement of the OHC vacuum condensate — a Hopf charge 
scar that cannot be smoothly removed.

The ratio between them is exactly $\alpha_0$ — the same 
number that determines the fine structure constant, the 
baryon asymmetry of the universe (Letter~71), and the 
quantum-classical boundary (Letter~72). The gravitational 
wave memory carries the fingerprint of the void geometry 
that underlies all of physics.

The Einstein Telescope~\cite{ET} and LISA~\cite{LISA} are 
currently under development. Standard GW memory searches are 
planned for both. BCT predicts that careful subtraction of 
the standard GR memory component will reveal a residual 
$\alpha_0 = 0.74\%$ topological component — the signature 
of the BCT vacuum lattice imprinted in spacetime.

\section{Conclusion}

The BCT Superfluid Lattice Model is fully consistent with 
all current gravitational wave observations. It makes two 
additional predictions: a scalar breathing mode at 
$h_\mathrm{scalar} = 1.61\times10^{-4}\,h_\mathrm{tensor}$, 
and a topological memory at $h_\mathrm{BCT} = \alpha_0 \times 
h_\mathrm{standard} = 1.30\times10^{-25}$ for GW150914-class 
events. Both are within the design sensitivity of the 
Einstein Telescope. The topological memory fraction $\alpha_0 
= r_\mathrm{oct}\times r_\mathrm{tet}/\pi$ is the void geometry 
signature of the BCT vacuum.

\begin{acknowledgments}
This letter was written in Melbourne, Australia, 17 March 2026, 
over the last of the birdseed money, while waiting for a car repair.
\end{acknowledgments}

\begin{thebibliography}{99}

\bibitem{GW150914}
B.~P. Abbott et al. (LIGO Scientific and Virgo Collaborations),
\textit{Phys. Rev. Lett.} \textbf{116}, 061102 (2016).

\bibitem{GWTC3}
R.~Abbott et al. (LIGO-Virgo-KAGRA),
\textit{Phys. Rev. X} \textbf{13}, 041039 (2023).

\bibitem{GW250114}
LIGO-Virgo-KAGRA Collaboration,
``Black Hole Spectroscopy and Tests of GR with GW250114,''
\textit{Phys. Rev. Lett.} (2026).

\bibitem{GW170817}
B.~P. Abbott et al.,
\textit{Phys. Rev. Lett.} \textbf{119}, 161101 (2017).

\bibitem{BCT_Monograph}
M.~R. Cabri\'{e},
``The BCT Superfluid Lattice Model,''
\href{https://doi.org/10.5281/zenodo.18884976}{doi:10.5281/zenodo.18884976} (2026).

\bibitem{BCT_AppJH}
M.~R. Cabri\'{e},
``BCT Appendix JH: The Hopfion Interior and the OHC,''
\href{https://doi.org/10.5281/zenodo.18975018}{doi:10.5281/zenodo.18975018} (2026).

\bibitem{GWMemory}
Y.~B. Zel'dovich and A.~G. Polnarev,
\textit{Sov. Astron.} \textbf{18}, 17 (1974).

\bibitem{ET}
M.~Punturo et al.,
\textit{Class. Quantum Grav.} \textbf{27}, 194002 (2010).

\bibitem{LISA}
P.~Amaro-Seoane et al.,
\textit{arXiv:1702.00786} (2017).

\end{thebibliography}

\end{document}
